% =====================================================================
\section{Inferring limit-order flow from the visible ladder}
\label{app:ladder}
% =====================================================================

The feed publishes the ten best price levels per side after every event, but not
the orders themselves. Submission and cancellation volume must therefore be
inferred from how the displayed depth changes. This appendix states the
accounting, the two rules that make it correct, what it cannot see, and the test
that checks it.

\subsection{The accounting identity}

Between two consecutive snapshots, the volume resting at a given price can change
only by execution, cancellation, or new submission:
%
\begin{equation}
V_n(p) \;=\; V_{n-1}(p) \;-\; \mathit{executed}_n(p) \;-\; \mathit{cancelled}_n(p)
\;+\; \mathit{submitted}_n(p).
\label{eq:ladder}
\end{equation}
%
Writing $\delta = V_n(p) - V_{n-1}(p)$ and letting $X_n(p)$ be the volume printed
at $p$ in that interval, equation~\eqref{eq:ladder} rearranges to
%
\begin{equation}
\mathit{submitted}_n(p) = \max\{\delta + X_n(p),\,0\},
\qquad
\mathit{cancelled}_n(p) = \max\{-(\delta + X_n(p)),\,0\}.
\end{equation}
%
Netting within the interval means an interval containing both a submission and a
cancellation at the same price reports only the difference. This biases both
flows downward, symmetrically, and the share statistics built from them are
therefore more robust than their levels.

\subsection{Two rules that make it correct}

\paragraph{Match by price, never by level index.}
When the book shifts a level --- the best ask is consumed and everything moves up
one slot --- an index-based comparison sees every level change and manufactures a
submission and a cancellation at each. Matching by price sees what actually
happened: one price emptied. This is the single largest source of error in
ladder-based order-flow inference, and the test in
Section~\ref{app:ladder-test} exercises it directly.

\paragraph{Count a price only where it is observable in both snapshots.}
On the ask side, every price at or below the deepest visible ask is observable:
if it is listed, its volume is shown, and if it is not listed, it is genuinely
empty, since a populated price inside the window would have displaced one of the
ten. Above that frontier nothing can be said. The comparison is therefore
restricted to
%
\[
p \;\le\; \min\!\left\{P^{10}_{n-1},\; P^{10}_{n}\right\}
\]
%
on the ask side, and the mirror condition on the bid side. Without this rule, a
price entering the window from beyond level ten arrives carrying pre-existing
depth that would be recorded as an enormous submission.

Two further restrictions are applied for the same reason. An interval is used
only when both of its endpoints show an ordinary two-sided quote, so nothing is
inferred across a special quote or a halt, when the displayed book does not mean
what it usually means; and only when the two rows are adjacent on the tape and in
the same session, so that a dropped stretch is never bridged. One unusable
snapshot therefore invalidates the two intervals that touch it. That is
deliberately conservative: it discards real flow rather than inventing any.

\subsection{What this cannot see}

Three limits are worth stating plainly, because they bound how the results in
Section~\ref{sec:book} should be read.

\begin{enumerate}[nosep]
\item \textbf{Depth beyond ten levels is invisible.} On the stock-days examined
      in detail, between 95\% and 98\% of quoted volume sits outside the ten
      displayed levels, in the aggregated \textsc{over} and \textsc{under}
      buckets. Ten levels of a 1-yen grid around a 4{,}000-yen price span a
      quarter of one percent, so this is arithmetic rather than a defect --- but
      it means the inferred flow describes the neighbourhood of the best quote,
      not the whole book. This matters most for the furthest distance class,
      which is precisely where Ohta argues the interesting orders sit, and it is
      the main reason the measures here are treated as complements to his rather
      than replacements.
\item \textbf{Order modifications are indistinguishable from a cancel-and-submit
      pair} at two different prices. For the digit shares used here this is
      harmless, since both the cancellation and the submission are attributed to
      the prices they actually occurred at.
\item \textbf{Off-book trading is excluded entirely.} Only the central limit
      order book is observed.
\end{enumerate}

One point cuts the other way and is worth recording: the TSE central book has no
hidden or iceberg orders, so displayed depth is the whole of the resting depth at
each visible price. The inference is exact where it applies.

\subsection{The test}
\label{app:ladder-test}

Correctness is checked against a hand-built order book whose flow is known
exactly, rather than against another implementation. The tape is six snapshots
long on a 1-yen grid with the ask side starting at 1{,}000, 1{,}001 and 1{,}002,
and it is constructed so that each of the three known failure modes would produce
a visibly wrong answer:

\begin{enumerate}[nosep]
\item 300 shares are added at 1{,}000 --- a submission, at the round price.
\item A buyer lifts 500 shares at 1{,}000. Displayed depth falls by 500, but the
      print explains all of it, so neither a submission nor a cancellation may be
      inferred. An implementation that forgets to net the execution reports 500
      shares of phantom cancellation.
\item 200 shares vanish at 1{,}000 with no print --- a cancellation.
\item The remaining 100 shares leave, 1{,}000 empties, the ladder shifts up, and
      1{,}003 enters from beyond the old window. The 100 must be counted; the
      depth arriving at 1{,}003 must not, and index matching would additionally
      report spurious flow at all three levels.
\end{enumerate}

The test asserts the exact expected totals --- 300 submitted, 300 cancelled, both
entirely at the round price, an execution ratio of $500/300$, and no flow at all
on the untouched bid side --- together with a separate case confirming that a
special quote in the middle of the tape invalidates the two intervals touching it
rather than being bridged across.

\subsection{Does it recover the published magnitudes?}

The strongest available check is external. Ohta reports mean $L^{S0}$ of 0.106
and mean $L^{S0C}$ of 0.806 across 2010--2022. Run on stock-days from 2024, this
implementation returns $L^{S0}$ between 0.103 and 0.155 and $L^{S0C}$ between
0.58 and 0.83. Nothing in the algorithm was calibrated against those numbers, and
they come from a different sample period, so the agreement is evidence that the
inference is recovering the same quantity the paper measures.
